\chapter{光滑性}\label{chap:smoothness}

\noindent\textbf{第 \thechapter{} 章学习目标}
\begin{itemize}
  \item 通过线性化, 直观地定义仿射簇的切空间.
  \item 在例子中计算仿射簇在一点处的线性化.
  \item 定义并使用向量空间的对偶空间.
  \item 用坐标环从内在角度刻画切空间.
  \item 描述仿射簇的维数与其切空间维数之间的关系.
  \item 判定仿射簇的光滑点和奇异点.
\end{itemize}

本书所画的仿射簇图形很可能让你想起了学微积分的日子. 你大概记得在微积分课上计算导数和切空间, 也可能记得导数和切空间可以用来刻画图像的\textquotedbl{}奇异性\textquotedbl{}出现在哪里. 例如, 绝对值函数 $f(x)=|x|$ 在 $x=0$ 处的导数没有定义, 对应于图像在 $x=0$ 的上方有一个\textquotedbl{}尖角\textquotedbl{}这一几何观察. 本章的目标是引入代数几何中的\textquotedbl{}微积分\textquotedbl{}. 具体而言, 要发展线性化和切空间的概念, 并用它们赋予仿射簇的奇异点以精确的含义.

本章先在第~\ref{sec:linearizations-and-tangent-spaces}~节定义切空间. 给定仿射簇 $X\subseteq\mathbb{A}^{n}$ 及一点 $a\in X$, 切空间 $T_{a}X$ 是 $K^{n}$ 的一个向量子空间, 其元素在几何上可视为在 $a$ 处与 $X$ 相切的向量. 但这一视角是外在的: 它有赖于把 $X$ 看成位于某个特定的仿射空间中. 我们在第~\ref{sec:vector-space-duality} 和~\ref{sec:intrinsic-tangent-spaces}~节处理这个问题. 我们先引入向量空间对偶的关键概念, 再用它们从坐标环 $K[X]$ 出发对切空间作出内在解释. 特别地, 这蕴含同构的仿射簇在对应点处的切空间同构.

作为向量空间, 切空间 $T_{a}X$ 的维数总是以 $X$ 的维数为下界, 这将在第~\ref{sec:dimension-of-tangent-spaces}~节证明. 切空间的维数与簇的维数相同的点称为簇的\emph{光滑点}, 而切空间的维数大于簇的维数的那些特殊点称为\emph{奇异点}. 本章最后在第~\ref{sec:smooth-and-singular-points}~节讨论光滑点和奇异点的性质.

\par\smallskip
{\footnotesize
\noindent\textcopyright{} The Author(s) 2025\par
\noindent E. Clader, D. Ross, \emph{Beginning in Algebraic Geometry}, Undergraduate Texts in Mathematics,
\href{https://doi.org/10.1007/978-3-031-88819-9_8}{\nolinkurl{https://doi.org/10.1007/978-3-031-88819-9_8}}.\par}
